% Lösungen Einheit 2 von 2 – Zinseszins und Guthabentabelle (zusammenbau v0.1)
\einheitenkopf{Einheit 2 von 2 · Zinseszins und Guthabentabelle}

%% TODO Lösung der Erklärzeile 15g) fehlt in der Bank
\erg{15}{a) vom neuen Guthaben 612\,€ \quad b) 927\,€ (Zinsen 27\,€) \quad c) 1\,632\,€ (Zinsen 32\,€) \quad d) 2\,310\,€ (Zinsen 110\,€) \quad e) 364\,€ (Zinsen 14\,€) \quad f) 5\,050\,€ (Zinsen 50\,€) \quad g) \ldots \quad h) 2\,121,80\,€ (nach einem Jahr 2\,060,00\,€, Zinsen im zweiten Jahr 61,80\,€) \quad i) Zinsen = 6,00\,€, Guthaben = 606,00\,€ \quad j) Guthaben = 477,00\,€, Zinsen = 9,54\,€ (Kontrolle: $477{,}00 + 9{,}54 + 120 = 606{,}54$) \quad k) 780\,€ ($750 \cdot 1{,}04$) \quad l) ≈ 3\,939,28\,€ ($3\,500 \cdot 1{,}03^{4}$) \quad m) Unterschied ≈ 10,91\,€ (einfach 4\,360\,€, Zinseszins ≈ 4\,370,91\,€) \quad n) ≈ 312,24\,€ (Endkapital ≈ 3\,312,24\,€) \quad o) 5\,000\,€ ($5\,306{,}04 : 1{,}02^{3}$) \quad p) Guthaben = 1\,081,82\,€, Zinsen = 32,45\,€ (Kontrolle: $1\,081{,}82 + 32{,}45 + 350 = 1\,464{,}27$)}
\erg{16}{a) 1\,003,97\,€ (Guthaben: 732,00\,€; 866,64\,€; 1\,003,97\,€)}
\erg{17}{a) 5\,955,08\,€ ($5\,000 \cdot 1{,}06^{3}$)}
\erg{18}{a) Die Zinsen kommen jedes Jahr vom neuen Guthaben, nicht vom Startkapital: richtig ≈ 2\,546,90\,€ ($2\,400 \cdot 1{,}02^{3}$). \quad b) Die 3\,\% werden jedes Jahr vom größeren Guthaben genommen – derselbe Anteil von mehr Geld ist mehr Geld. \quad c) Ja, knapp: nach 4 Jahren ≈ 2\,251,02\,€. \quad d) 4\,\% (Wachstumsfaktor $2\,080 : 2\,000 = 1{,}04$)}
